\documentclass[asi,article,accept,pdftex,moreauthors]{Definitions/mdpi}
\usepackage{amsthm}
\usepackage{adjustbox}
\usepackage{booktabs}
\usepackage{amsmath, amssymb}
\usepackage{xcolor}
\usepackage{graphicx}
\firstpage{1} 
\makeatletter 
\setcounter{page}{\@firstpage} 
\makeatother
\pubvolume{9}
\issuenum{1}
\articlenumber{0}
\pubyear{2026}
\copyrightyear{2026}
\externaleditor{Firstname Lastname} % More than 1 editor, please add `` and '' before the last editor name
\datereceived{8 October 2026} 
\daterevised{12 November 2026} % Comment out if no revised date
\dateaccepted{18 November 2026} 
\datepublished{ } 
%\newtheorem{proposition}{Proposition}
%\newtheorem{remark}{Remark}

% ==== Round-2 revision markup =======================================
% \rev{...} marks the NEW corrections of this (Round 2) iteration -> BLUE.
% Round 1 text (written with \textcolor{blue}{...}) is forced to BLACK here,
% so only this round's changes appear in color.
%\let\oldtextcolor\textcolor
\newcommand{\rev}[1]{#1}% NORMAL version: Round-2 additions in black
\let\oldtextcolor\textcolor
\renewcommand{\textcolor}[2]{\ifx\relax#1\relax#2\else\oldtextcolor{#1}{#2}\fi}% neutralize Round-1 \textcolor{blue}{...} to black
% ====================================================================

\Title{Koopman-Based State Observer for a Coupled \linebreak  Three-Tank System}

\Author{
  Wilder Hernandez Manosalva,
  Carlos Andres Torres-Pinzon * and 
  Duvan Tellez-Castro
}
\AuthorNames{Wilder Steven Hernandez Manosalva, Carlos Andrés Torres-Pinzón, Duván Téllez-Castro}

\address[1]{
Facultad de Ingeniería, Universidad Distrital Francisco José de Caldas, Bogotá 110231, Colombia; caatorresp@udistrital.edu.co (C.A.T.-P.); duatellezc@udistrital.edu.co (D.T.-C.)}
\corres{Correspondence: caatorresp@udistrital.edu.co}

%-----------------------------------------------------------------
\abstract{%
This paper presents a Koopman-based Luenberger-type observer that reconstructs the unmeasured level of a partially observed three-tank hydraulic system. {The observer operates natively in an eight-dimensional space of physics-informed observables identified via Extended Dynamic Mode Decomposition (EDMD). Its feedback gain is synthesized offline through the discrete algebraic Riccati equation (DARE), and the ultimate boundedness of the estimation error is established analytically in a weighted Lyapunov norm.} The observer is benchmarked in a closed loop against an Extended Kalman Filter (EKF) under an identical model predictive controller, over five independent runs per observer (ten in total) on a physical plant instrumented with an independent ground-truth level sensor on the unmeasured tank. {Under the as-run experimental gains, the Koopman observer reduces the observation RMSE against ground truth by 19.3\% ($1.139 \pm 0.057$ vs.\ $1.412 \pm 0.064$~cm), whereas the physical tracking error of the plant against the setpoint shows no statistically detectable difference ($0.804 \pm 0.031$ vs.\ $0.799 \pm 0.011$~cm, Welch $p=0.74$). An offline leave-one-run-out cross-validated replay of the recorded telemetry indicates that this initial observation advantage is largely attributable to the specific EKF noise-covariance ratios utilized in the live experiments. When both observers are tuned via cross-validation, they achieve comparable estimation accuracy ($0.655 \pm 0.023$~cm for Koopman vs.\ $0.600 \pm 0.020$~cm for EKF). A predictor-only baseline confirms that the output injection actively rejects a deliberately injected $+2.0$~cm offset, decaying it to $10\%$ within approximately $7$~s in the replay ($12$~s under noise-free propagation), while the pure predictor drifts to $>5.5$~cm RMSE. Isolated algorithmic profiling confirms that the full recursion of an optimized EKF ($0.087$~ms with the AD Jacobian, against $0.038$~ms for the Jacobian evaluation alone) is comparable to the linear Koopman update ($0.025$~ms); both are vastly overshadowed by the non-convex predictive control optimization, which alone requires $14.50$~ms on average (P99 $19.90$~ms) and yields a complete Koopman observer-plus-controller pipeline of $14.53$~ms on average (P99 $19.96$~ms; the corresponding EKF pipeline is $15.03$~ms with the finite-difference Jacobian and $14.59$~ms with the algorithmic-differentiation Jacobian), negligible against the $2.1$~s hardware sampling period. The Koopman observer therefore delivers EKF-grade precision with a linear, time-invariant structure, eliminating real-time Jacobians while providing an ultimate-boundedness guarantee for the lifted estimation error of the identified model.}
}

\keyword{Extended Kalman Filter; Koopman operator; {nonlinear programming}; three-tank system; state estimation}


%\pubyear{2026}
%\copyrightyear{2026}

%% =====================================================================
\begin{document}
%% =====================================================================
\section{Introduction}\label{sec:intro}

The problem of state estimation for nonlinear systems based on partial observations remains a key issue in control engineering. In hydraulic processes such as coupled-tank systems~\cite{johansson2000}, level sensors are often impractical because of their cost, limited accessibility, or failure, so unmeasured states must be estimated from the available measurements. Such an observer should not only be accurate, but also fast and robust to measurement noise, which is especially important when it is part of a control loop with a fixed sampling rate.

The Extended Kalman Filter (EKF) is the observer method most frequently used for nonlinear systems. The Extended Kalman Filter uses a linearization process to update the state using the Jacobian Matrix~\cite{simon2006}. This method, although quite powerful in numerous applications, involves computational effort proportional to the complexity of the Jacobian. This may be a drawback, especially for stiff or large problems, where real-time implementation is difficult. Furthermore, the EKF may perform poorly for highly nonlinear processes in which the linearization assumption fails, or when the estimate is far from the true trajectory during reference changes~\cite{julier2004}.

In the current work, the EKF is used instead of the Moving Horizon Estimator (MHE) because both the EKF and the Koopman-Luenberger observer presented here are one-shot recursions where each recursion provides an estimate using only one prediction-correction cycle. In contrast, an MHE has to solve a nonlinear program with constraints in a sliding window at each sampling point, requiring a computational effort equivalent to or even greater than that of the MPC algorithm. With the same $T_s = 2$~s and the computational budget of a single-board computer, an MHE would not have been a like-for-like comparison. The EKF, in contrast, is the standard baseline for benchmarking nonlinear estimators. An MHE nevertheless remains a valid alternative, since it exploits information from a sliding window rather than a single correction step~\cite{rao2003}.


As an alternative to linearization of the nonlinear dynamics locally at each time step, we employ a linear approximation of the dynamics of the system over a set of observables in a higher dimensional space, for example, through Extended Dynamic Mode Decomposition (EDMD)~\cite{williams2015,korda2018}. Here, the Koopman operator~\cite{koopman1931, mezic2005, brunton2016} serves as a linear infinite-dimensional representation of the nonlinear flow and can be employed for finite-dimensional reduction and spectral decomposition of the system. In this setting, classical linear techniques such as Luenberger observers and MPC can be applied to the original nonlinear system without repeated Jacobian computations. The approach is especially promising for hydraulic processes, since their main nonlinearity, the Torricelli discharge law, can be described compactly with physically motivated observables.


Koopman-based control and estimation have been explored in several settings~\cite{tellez2022,hernandez2026,aponte2025,DURANSIGUENZA2025,netto2018}. Surana and Banaszuk formulated linear observers in lifted spaces~\cite{surana2016}; later work addressed constrained estimation~\cite{Surana2017} and hybrid Koopman--EKF schemes~\cite{kiloekf2026}. {The Luenberger-type structure in lifted coordinates is therefore established in the literature and is not claimed as a theoretical novelty here. What remains scarce is a rigorous, closed-loop experimental comparison with an EKF on a physical plant incorporating an independent ground-truth sensor on the unmeasured state, combined with an offline cross-validated replay that decouples observer accuracy from closed-loop interaction and tuning artifacts.}

{The primary contributions of this paper are three-fold. First, in terms of observer design, we construct a physics-informed dictionary of eight observables encoding Torricelli square-root discharge and pairwise level couplings, synthesizing a time-invariant gain offline via the DARE in the lifted space. Second, regarding error analysis, we provide an ultimate-boundedness analysis for the lifted estimation error. Although it builds on standard linear-systems theory, its contribution is the specialization to the lifted EDMD setting with explicit residual decomposition under EDMD residuals and measurement noise, evaluated in a weighted Lyapunov norm. Third, as an experimental benchmark, we present ten closed-loop physical runs utilizing an independent ground-truth sensor. This benchmark uses three error signals to separate estimation accuracy from physical tracking error, together with a unified steady-state window and an offline leave-one-run-out cross-validation replay that supports a fair comparison.}

\rev{Relative to existing Koopman-based state-estimation approaches~\cite{surana2016,Surana2017,kiloekf2026}, the practical novelty of this work is not the lifted-space observer synthesis itself but the joint demonstration, on a single physical plant, of three elements that are seldom reported together: (i) a gain designed offline through the DARE with an a~priori computable ultimate-bound certificate for the lifted error; (ii) a direct separation of the observation error $\hat h_2-h_2^\mathrm{real}$ from the closed-loop control residual $\hat h_2-h_2^\mathrm{ref}$ by means of an independent ground-truth sensor on the unmeasured tank; and (iii) a like-for-like closed-loop comparison against the EKF under an identical predictive controller, complemented by an offline cross-validated replay that disentangles the estimator structure from its tuning. In the experimental Koopman-observer literature, validation is typically indirect (against the model or the controller target) rather than against an independent physical measurement of the unobserved state.}

%% =====================================================================
\section{System Description}
\label{sec:system}

\subsection{Physical Plant}

The experimental platform consists of three identical cylindrical tanks fabricated
from transparent PVC pipe (inner diameter $6.2$~cm, cross-sectional area
$A_\mathrm{tank} = 3.02 \times 10^{-3}$~m$^2$), arranged horizontally and coupled
by needle valves. The tanks feed into a common reservoir from which two R385 diaphragm
pumps recirculate water back to tanks~1 and~3. The coupling valves between tanks~1--2 and tanks~2--3 allow bidirectional flow governed by the head difference between adjacent tanks. The choice of this configuration---with control inputs to tanks~1 and~3 but no direct sensor on tank~2---deliberately creates a partially observable system where $h_2$ must be reconstructed dynamically.

\subsection{Ground-Truth Validation Sensor and Data Acquisition}
\label{sec:ground_truth}

{To validate the estimated level $\hat{h}_2$ against physical reality, an independent HC-SR04 ultrasonic sensor was installed above tank~2. This sensor is used exclusively for post-hoc evaluation; its readings (\rev{$h_2^\mathrm{real}$}) are logged alongside each run but are completely isolated from both the state observers and the predictive controller. Tank~2 therefore remains strictly unmeasured from the perspective of the closed control loop.}

{The ESP32 microcontroller firmware dictates the data acquisition timing, which is crucial for interpreting the physical results. Each fluid level is obtained as the median of seven consecutive ultrasonic pulses. The ESP32 aggregates these readings and transmits them to the host PC only after a $2000$~ms interval has elapsed, checked at the end of each acquisition cycle. As a result, the effective sampling period seen by the host is $T_s^\mathrm{eff} \approx 2.10$~s (P99 $2.15$~s), whereas the model and the controller assume a nominal $T_s = 2.0$~s. Actuation commands are also delayed in transmission: the pump voltages are applied approximately $0.7$~s after the host issues the command. In steady state, the ground-truth readings exhibit an empirical standard deviation of $0.45$~cm (ranging between $0.35$ and $0.58$~cm). This variability reflects real level fluctuations caused by the recirculating pump jets, rather than purely electrical transducer noise, in addition to the $3$~mm quantization resolution of the HC-SR04.}

\subsection{Mathematical Model}

The three-tank dynamics follow mass conservation with Torricelli discharge and
viscous-turbulent transitional coupling flows:

\begin{equation}
\frac{d}{dt}\begin{pmatrix}h_1\\h_2\\h_3\end{pmatrix}
= \frac{1}{A_\mathrm{tank}}
\begin{pmatrix}
  q_{i1} - q_{12} - q_{o1}\\
  q_{12} - q_{23} - q_{o2}\\
  q_{i3} + q_{23} - q_{o3}
\end{pmatrix},
\label{eq:ode}
\end{equation}
where $q_{oi} = \alpha_{oi} A_{oi} \sqrt{2g h_i}$ represents the Torricelli free discharge through the orifice area $A_{oi} = \pi D_{oi}^2 / 4$. The inter-tank flows $q_{12}$ and $q_{23}$ capture transitional pipe friction:

\begin{equation}
q_{12} = \alpha_{12}(h_1,h_2)\,A_{12}\sqrt{2g|h_1-h_2|}\,\operatorname{sgn}(h_1-h_2),
\quad
\alpha_{12} = \alpha_{120}\tanh\!\left(\!\frac{2\lambda_{12}}{\lambda_{c12}}\!\right),
\end{equation}
with the flow parameter $\lambda_{12} = D_{12}\rho\eta^{-1}\sqrt{2g|h_1-h_2|}$, and analogously for $q_{23}$. The physical parameters, calibrated from tank-drain experiments and open-loop step tests, are listed in Table~\ref{tab:params}.

\begin{table}[H]
 % \centering
  \caption{Calibrated physical parameters of the three-tank experimental plant.}
  \label{tab:params}
  \setlength{\tabcolsep}{6.1mm}{
  \begin{tabular}{llll}
    \toprule
    \textbf{Parameter} & \textbf{Value} & \textbf{Unit} & \textbf{Description} \\
    \midrule
    $A_\mathrm{tank}$ & $3.02\times10^{-3}$ & m$^2$        & Cross-sectional area \\
    $\rho$            & 997                 & kg/m$^3$     & Water density \\
    $\eta$            & $8.9\times10^{-4}$  & Pa$\cdot$s   & Dynamic viscosity \\
    $\alpha_{o1,2,3}$ & 0.0271              & ---          & Discharge coefficient \\
    $D_{o1,3}$        & 10                  & mm           & Outlet orifice diameter \\
    $\alpha_{120}$    & 0.3038              & ---          & Turbulent coupling coeff.\ 1--2 \\
    $\alpha_{230}$    & 0.1344              & ---          & Turbulent coupling coeff.\ 2--3 \\
    $D_{12}$, $D_{23}$& 10                  & mm           & Coupling orifice diameter \\
    $\lambda_{c12}$   & 24000               & ---          & Critical flow number, 1--2 \\
    $\lambda_{c23}$   & 29600               & ---          & Critical flow number, 2--3 \\
    $q_{i1}^\mathrm{max}$ & $2.69\times10^{-5}$ & m$^3$/s  & {Max.\ pump to tank~1 flow} \\
    $q_{i3}^\mathrm{max}$ & $2.41\times10^{-5}$ & m$^3$/s  & {Max.\ pump to tank~3 flow} \\
    $T_s$             & 2.0                 & s            & Model sampling period \\
    \bottomrule
  \end{tabular}}
\end{table}

%% =====================================================================
\section{Methods}
\label{sec:methods}

\subsection{EDMD Identification and Spectral Stabilization}
\label{sec:edmd}

A dictionary of 8 physics-informed observables was selected based on the Torricelli
nonlinearity and pairwise couplings in the three-tank model:

\begin{equation}
\Psi(x) = \bigl[\sqrt{h_1},\;\sqrt{h_2},\;\sqrt{h_3},\;h_1h_2,\;h_2h_3,\;h_1h_3,\;h_1,\;h_3\bigr]^\top \in \mathbb{R}^8.
\label{eq:psi}
\end{equation}

The square-root terms directly encode the Torricelli discharge nonlinearity, while the bilinear cross-terms $h_ih_j$ capture the state-dependent coupling coefficients.

{Training Data Generation.
 Because $h_2$ remains unmeasured in the physical control loop, the identification data utilized for EDMD was synthesized entirely offline. The calibrated continuous-time ODE~\eqref{eq:ode} was integrated via a 4th-order Runge-Kutta scheme at $T_s=2.0$~s to generate $M=8220$ dynamic snapshots. This dataset encompasses simulated closed-loop trajectories around multiple equilibria, large state perturbations ($\pm25$--$40\%$), and dynamic reference transitions that include the setpoint levels later used in the experiments. Generating training data via the mathematical simulator guarantees exact access to the true unmeasured state $h_2$ for dictionary construction, completely isolated from physical sensor noise.}

{The EDMD regression is formulated as a ridge-regularized least-squares optimization on the stacked regressors $[\Psi(x_k)^\top\;u_k^\top]$:}
\begin{equation}
{[\mathbf{A}_\mathrm{raw},\mathbf{B}] = \arg\min_{A,B} \sum_{k=1}^{M}
  \bigl\|\Psi(x_{k+1}) - \mathbf{A}\Psi(x_k) - \mathbf{B}u_k\bigr\|_2^2
  + \lambda_R\bigl(\|\mathbf{A}\|_F^2+\|\mathbf{B}\|_F^2\bigr), \quad \lambda_R=10^{-4}.}
\label{eq:edmd}
\end{equation}

{To obtain a Schur-stable predictor, a spectral projection was applied after identification: $\mathbf A=\mathbf V\operatorname{diag}(\tilde\lambda)\mathbf V^{-1}$, shrinking any eigenvalue $|\lambda_i| \ge 1$ to $1-\delta$ ($\delta=10^{-4}$) and leaving the remaining eigenvalues unchanged. The resulting state matrix has spectral radius $\rho(\mathbf{A}) = 0.99990 = 1-\delta$, attained after one eigenvalue $|\lambda_i| \ge 1$ was rescaled to $1-\delta$. When evaluated on ten independent held-out simulated trajectories ($M=2000$ snapshots, validation seed $123$ against the training seed $42$), the 20-step predictive RMSE for $h_2$ is $0.020$~cm, confirming the high fidelity of the lifted linear approximation relative to the underlying ODE simulator.}

\rev{Although the EDMD matrices are fitted on simulated trajectories, the simulator is not a generic model: its parameters (discharge and coupling coefficients, orifice areas, critical flow numbers in Table~\ref{tab:params}) were identified from physical experiments on the same plant (free-drain and single-pump step tests), so the training distribution is anchored to the measured hardware rather than to nominal values. Residual model--plant mismatch is not assumed away: it enters the observer error dynamics~\eqref{eq:error_dyn} as the bounded residual $r_k$, for which Proposition~\ref{prop:stability} provides the ultimate bound. The experimental results are the end-to-end evidence that this mismatch is tolerable: all ten closed-loop runs remain stable without any online re-identification, and the predictor-only baseline of Section~\ref{sec:predictor} shows that the output injection actively rejects the accumulated open-loop drift, keeping the observation RMSE bounded at $\approx 1.15$~cm with the as-run gain ($0.655$~cm after tuning, Section~\ref{sec:tuning_cv}), against $5.5$--$6.1$~cm for the pure predictor, rather than diverging. The remaining observation bias exceeds the $\approx 0.45$~cm sensor noise, which indicates that model--plant mismatch, not sensor noise, dominates the residual error; this is acknowledged in Section~\ref{sec:limitations}.}

\subsection{Koopman Observer Design and Stability}
\label{sec:koopman_obs}

\subsubsection{Observer Structure and Measurement Model}

The observer propagates the lifted state vector $\hat{z}_k \in \mathbb{R}^8$:
\begin{align}
\hat z_{k|k-1}&=\mathbf A\hat z_{k-1}+\mathbf Bu_{k-1},\label{eq:pred}\\
\hat z_k&=\hat z_{k|k-1}+L\bigl(y_k-C\hat z_{k|k-1}\bigr),\label{eq:corr}
\end{align}
where $y_k = [\sqrt{h_{1,\mathrm{meas}}},\;\sqrt{h_{3,\mathrm{meas}}},\;h_{1,\mathrm{meas}},\;h_{3,\mathrm{meas}}]^\top$
and $C \in \mathbb{R}^{4\times8}$ extracts the components 
$\sqrt{h_1}$, $\sqrt{h_3}$, $h_1$, and $h_3$ from $\Psi$ 
(i.e., entries 1, 3, 7, and 8 in 1-based indexing). {The measurement model is explicitly $y_k = C z_k + v_k$, where $C$ selects the observable components $[\sqrt{h_1},\,\sqrt{h_3},\,h_1,\,h_3]^\top$ and $v_k\in\mathbb R^4$ collects the corresponding measurement noise.} The pair $(\mathbf{A}, C)$ possesses full observability rank. The physical estimate is recovered via $\hat h_i=(\max(0,\hat z_i))^2$, where the non-negativity clipping operation prevents algorithmic failure from non-physical negative heights during aggressive transient maneuvers.

{Measurement Noise Propagation. Assuming the raw physical ultrasonic sensors are affected by zero-mean noise $\nu=[\nu_1,\nu_3]^\top\sim(0,\Sigma_\nu)$, a first-order Delta method expansion characterizes the noise on the augmented output vector as $v_k\approx \mathbf J_\nu(x)\nu_k$, yielding a transformed covariance $\Sigma_y(x) = \mathbf J_\nu\Sigma_\nu\mathbf J_\nu^\top$. Because the linear head variables and their respective square roots originate algebraically from the exact same physical transducers, the resulting $4 \times 4$ covariance matrix $\Sigma_y(x)$ is structurally singular (rank 2). Consequently, utilizing a strictly positive-definite measurement penalty $R \succ 0$ during Riccati synthesis represents a necessary control-theoretic loop-shaping regularization rather than a probabilistically exact maximum-likelihood covariance inversion.}

{The observer gain $L$ is synthesized by solving the discrete algebraic Riccati equation (DARE) in predictor form:
\begin{equation}
L=\mathbf AXC^\top(CXC^\top+R)^{-1},\qquad
X=\mathbf AX\mathbf A^\top-\mathbf AXC^\top(CXC^\top+R)^{-1}CX\mathbf A^\top+Q.
\label{eq:dare}
\end{equation}

During live physical experiments, the penalties were defined as $Q=I_8$ \mbox{(i.e., $\alpha=1$)} and $R=I_4$. The closed-loop error matrix governing the implemented observer is $\mathbf A_c=(I-LC)\mathbf A$, characterized by a spectral radius $\rho(\mathbf A_c)=0.93515$, a maximum gain magnitude $|L|_{\max}=0.5867$, and a decay time constant $\tau_e=-T_s/\ln\rho(\mathbf A_c)=29.83$~s. The numerical values of the identified matrices $\mathbf A$ and $\mathbf B$, the output selector $C$, and the observer gain $L$ are reported in Appendix~\ref{app:matrices}. An alternative tuning configuration, derived later via offline cross-validation (Section~\ref{sec:tuning_cv}), uses \rev{a diagonal $Q$ whose entries} associated with the measured components of the lifted state $[\sqrt{h_1},\sqrt{h_3},h_1,h_3]^\top$ are set to $\alpha_\mathrm{med}=10^{-3}$ and \rev{whose} remaining entries \rev{are set} to $\alpha_\mathrm{un}=0.1$.}

\subsubsection{Stability and Error Boundedness}
\label{sec:stability}

Let $e_k=z_k-\hat z_k$ denote the estimation error in the lifted coordinate space, where $z_k=\Psi(x_k)$. Accounting for the EDMD truncation residual $r_k=z_{k+1}-(\mathbf Az_k+\mathbf Bu_k)$ and the output noise $v_k$, the error dynamics obey:
\begin{equation}
e_{k} = \mathbf A_c\,e_{k-1}+\varepsilon_k,\qquad \mathbf A_c\triangleq(I-LC)\mathbf A,\qquad
\varepsilon_k\triangleq(I-LC)r_{k-1}-Lv_k .
\label{eq:error_dyn}
\end{equation}

\begin{Proposition}[Ultimate boundedness of the lifted estimation error]
\label{prop:stability}
{Assume $R\succ0$, $Q\succ0$, $(\mathbf A,C)$ observable, and $(\mathbf A,Q^{1/2})$ stabilizable. Let $\mathbf A_c=(I-LC)\mathbf A$ be Schur stable (verified numerically for the as-run gain, $\rho(\mathbf A_c)=0.93515$). Let $W\succ0$ denote the unique solution to the discrete Lyapunov equation $\mathbf A_c^\top W\mathbf A_c-W=-I$, defining the weighted norm $\|e\|_W=\sqrt{e^\top We}$ and contraction factor $\gamma=\sqrt{1-1/\lambda_{\max}(W)}<1$. Under uniformly bounded modeling residuals ($\|r_k\|_2\le\bar r$) and measurement noise ($\|v_k\|_2\le\bar v$), the estimation error obeys the ultimate bound:
\begin{equation*}
\limsup_{k\to\infty}\|e_k\|_2\le\frac{\|I-LC\|_2\,\bar r+\|L\|_2\,\bar v}{1-\gamma}\sqrt{\kappa(W)},
\end{equation*}
where $\kappa(W) = \lambda_{\max}(W)/\lambda_{\min}(W)$ represents the condition number of the Lyapunov metric.}
\end{Proposition}

\begin{proof}
{The strict positivity of $Q$ and $R$ alongside observability of $(\mathbf A,C)$ guarantees the existence of a stabilizing DARE solution. In the present implementation the DARE is solved in predictor form, $L=\mathbf AXC^\top(CXC^\top+R)^{-1}$, which stabilizes $\mathbf A-LC$. Because the implemented correction step~\eqref{eq:corr} uses the same $L$ as a filter gain, the closed-loop matrix governing~\eqref{eq:error_dyn} is $\mathbf A_c=(I-LC)\mathbf A$, which differs from $\mathbf A-LC$; the DARE therefore provides a design procedure for $L$ but not a direct stability guarantee for $\mathbf A_c$. Schur stability of $\mathbf A_c$ is verified numerically for the as-run gain ($\rho(\mathbf A_c)=0.93515$), and $W$ denotes the corresponding unique solution of the discrete Lyapunov equation $\mathbf A_c^\top W\mathbf A_c-W=-I$. From this identity the induced $W$-norm satisfies $\gamma^2=1-1/\lambda_{\max}(W)<1$. The triangle inequality yields $\|e_k\|_W \le \gamma\|e_{k-1}\|_W+\|\varepsilon_k\|_W$. Bounding the disturbance norm as $\|\varepsilon_k\|_W \le \sqrt{\lambda_{\max}(W)}\,(\|I-LC\|_2\bar r+\|L\|_2\bar v)$ and taking the limit superior as $k \to \infty$ establishes the asymptotic boundary. The standard Euclidean bound is subsequently recovered via the inequality $\|e\|_2 \le \|e\|_W / \sqrt{\lambda_{\min}(W)}$.}
\end{proof}

\begin{Remark}
{For the experimental observer utilizing $\alpha=1$, the Lyapunov equation is solved for the same closed-loop matrix $\mathbf A_c=(I-LC)\mathbf A$ that governs the error dynamics~\eqref{eq:error_dyn}, yielding $\lambda_{\max}(W)=11.33$, $\lambda_{\min}(W)=1.001$, $\kappa(W)=11.31$, and amplification factor $\sqrt{\kappa(W)}/(1-\gamma)=74.5$, confirming a strict contraction factor $\gamma=0.9548 < 1$. Because the theoretical maximum bound accommodates the worst-case spectral alignment of the EDMD truncation error $\bar r$ and ultrasonic sensor noise $\bar v$ across all operating regimes, the resulting algebraic threshold is highly conservative. {The bound concerns the identified lifted model; on the real plant, the effective residual $\varepsilon_k$ includes the multi-step EDMD error and is bounded only on the region covered by the training data.} \rev{To make the design-versus-implementation relationship explicit: the DARE~\eqref{eq:dare} yields a gain that stabilizes the \emph{predictor} matrix $\mathbf A-LC$, and is therefore used here only as a gain-\emph{design} procedure; the matrix that actually governs the sample-to-sample propagation of the implemented error~\eqref{eq:error_dyn} is the filter-form matrix $\mathbf A_c=(I-LC)\mathbf A$, whose Schur stability is verified numerically for the as-run gain ($\rho(\mathbf A_c)=0.93515$) and certified by the Lyapunov solution $W$ above. All stability quantities in this paper ($\gamma$, $\kappa(W)$, $\tau_e$) are computed from $\mathbf A_c$, not from $\mathbf A-LC$.} Proposition~\ref{prop:stability} formally guarantees mathematical boundedness without serving as an empirical predictor of tight operational accuracy.}
\end{Remark}

\subsection{Extended Kalman Filter Baseline}
\label{sec:ekf}

The EKF propagates the physical state $\hat{x}_k = [\hat{h}_1,\hat{h}_2,\hat{h}_3]^\top$ via the continuous-time ODE~\eqref{eq:ode} integrated with RK4. {To ensure a rigorous baseline, two implementations of the continuous-time Jacobian $J_k=\partial f/\partial x$ were used. During the live experiments, the Jacobian was evaluated by forward finite differences, as is common in practical implementations. Subsequently, during the offline algorithmic profiling phase (Section~\ref{sec:timing}), an exact Jacobian generated with CasADi algorithmic differentiation (AD) was used as an efficient reference implementation. All EKF accuracy results in this paper correspond to the finite-difference implementation deployed in the live runs.} The prediction and correction steps follow standard Gaussian propagation formulas~\cite{simon2006}. The covariance matrices deployed during live runs were $Q_\mathrm{EKF} = \mathrm{diag}[(0.5\,\mathrm{mm})^2,\,(2\,\mathrm{mm})^2,\,(0.5\,\mathrm{mm})^2]$ and $R_\mathrm{EKF} = \mathrm{diag}[(1.5\,\mathrm{mm})^2,\,(1.5\,\mathrm{mm})^2]$.

\subsection{Model Predictive Controller (Nonlinear Program)}
\label{sec:mpc}

Both observers feed the identical model predictive controller (MPC), so that performance differences can be attributed to the state estimator. {The MPC is formulated over the lifted EDMD model $(\mathbf{A},\mathbf{B})$. Because the primary controlled variable is the physical liquid level $h_2$, which is represented algebraically in the lifted space as $h_2 = z_2^2$, tracking the reference trajectory naturally imposes a quartic cost function:}
\begin{equation}
{\min_{\{\tilde{u}_n\}_{n=0}^{N-1}}
\sum_{k=1}^{N} w_k\,q\,\bigl(\tilde{z}_{k,2}^2 - h_2^\mathrm{ref}\bigr)^2
+ \sum_{n=0}^{N-1}\!\left(\tilde{u}_n^\top R_1\tilde{u}_n + \Delta\tilde{u}_n^\top R_2\Delta\tilde{u}_n\right).}
\label{eq:mpc}
\end{equation}

{Subject to the linear equality constraints $\tilde{z}_{n+1} = \mathbf{A}\tilde{z}_n + \mathbf{B}\tilde{u}_n$ and actuation bounds $0 \le \tilde{u}_n \le 1$, the optimization problem constitutes a non-convex Nonlinear Program (NLP). This NLP is solved at each sampling instance by the interior-point optimizer IPOPT~\cite{wachter2006} via the CasADi interface~\cite{andersson2019}; IPOPT factorizes the sparse symmetric indefinite KKT system using a sparse direct linear solver (MUMPS in the CasADi-bundled IPOPT build; HSL MA27/MA57 if installed).} Design parameters are defined as $N = 20$, $q = 500$, $R_1 = \mathrm{diag}(0.1,\,1.0)$, and $R_2 = \mathrm{diag}(0.1,\,0.1)$, with stage weight $w_k=1$ for $k<N$ and terminal weight $w_N=10$. The optimizer is warm-started at each step from the previous solution shifted by one sample: {$s_0^{(k)}=[\tilde z_{1:N}^{(k-1)};\tilde z_N^{(k-1)};\tilde u_{1:N-1}^{(k-1)};\tilde u_{N-1}^{(k-1)}]$}. The heavier penalty on pump~3 reflects its lower maximum flow capacity. When the EKF feeds the MPC, its physical state vector is simply lifted before optimization: $\hat{z}_k = \Psi(\hat{x}_{k|k})$ (Figure~\ref{fig:architecture}).

\begin{figure}[H]
%  \centering
  \includegraphics[width=\linewidth]{img/architecture_diagram.png}
  \caption{Closed-loop control architecture for both pipelines. \textbf{Left:} Koopman Observer pipeline evolving natively in the lifted space $\mathbb{R}^8$. \textbf{Right:} EKF pipeline estimating in $\mathbb{R}^3$ and lifting via $\Psi(\cdot)$ prior to optimization. Both pipelines share identical MPC parameters ($N=20$, $q = 500$), references, and physical plant connections.}
  \label{fig:architecture}
\end{figure}

\subsection{Performance Metrics and Unified Evaluation Window}
\label{sec:protocol}

Five independent closed-loop runs of 400 samples each (nominal duration 800~s) were executed sequentially on the physical plant for each observer. The same setpoint profile, with one ascending and one descending step, was used in all runs:
\begin{equation}
h_2^\mathrm{ref}(t) = \begin{cases}
  0.10\;\mathrm{m} & 0 \le t < 264\;\mathrm{s}, \\
  0.20\;\mathrm{m} & 264 \le t < 528\;\mathrm{s}, \\
  0.15\;\mathrm{m} & 528 \le t \le 800\;\mathrm{s}.
\end{cases}
\label{eq:ref}
\end{equation}

{The reference is indexed by sample: it switches at $k=132$ and $k=264$ ($t_k=264$ and $528$~s nominal). To separate estimation accuracy from closed-loop tracking behavior, three error metrics are evaluated at every sample. The controller tracking error, defined as $e^\mathrm{track}_k = \hat h_{2,k} - h^\mathrm{ref}_{2,k}$, measures how well the MPC regulates the internal state estimate to the setpoint. The observation error, defined as $e^\mathrm{obs}_k = \hat h_{2,k} - h^\mathrm{real}_{2,k}$, benchmarks the accuracy of the estimator strictly against the independent ground-truth sensor. Finally, the physical tracking error, defined as $e^\mathrm{real}_k = h^\mathrm{real}_{2,k} - h^\mathrm{ref}_{2,k}$, quantifies the true physical position of the fluid level relative to the desired setpoint. These quantities satisfy the algebraic identity $e^\mathrm{real}_k = e^\mathrm{track}_k - e^\mathrm{obs}_k$, and the superscripts (track / obs / real) are used consistently in every table and figure caption that follows. \rev{For notational uniformity throughout the paper, the lifted state is written $z$ (with estimate $\hat z$ and error $e=z-\hat z$), the EDMD operators are set in bold ($\mathbf A$, $\mathbf B$) with the output selector $C$, the DARE solution is $X$ and the Lyapunov certificate is $W$, and the implemented closed-loop error matrix is denoted $\mathbf A_c=(I-LC)\mathbf A$ wherever it appears. The ground-truth level is written $h_2^\mathrm{real}$ throughout, $T_s$ denotes only the sampling period, and settling times are denoted $t_s$.} All statistical metrics employ a single unified steady-state window. Exactly $60$~s of transient data following each reference transition (intervals $[0,60)$, $[264,324)$, and $[528,588)$~s, i.e.,\ $3\times 30 = 90$ samples excluded out of 400) are removed. The retained steady-state intervals $[60,264)$, $[324,528)$ and $[588,800)$ contain $102$, $102$ and $106$ samples respectively. In discrete sample indices $k\in\{0,\ldots,399\}$ with $t_k=2k$~s, the retained indices are $k\in\{30,\ldots,131\}\cup\{162,\ldots,263\}\cup\{294,\ldots,399\}$, yielding precisely $N_\mathrm{SS} = 102+102+106 = 310$ steady-state samples per run. All temporal axes in the manuscript are expressed in nominal time ($T_s=2.0$~s); the physical acquisition chain operates at $T_s^\mathrm{eff}\approx 2.10$~s (Section~\ref{sec:ground_truth}), so a nominal $60$~s window corresponds to approximately $63$~s of wall-clock time.}
%% =====================================================================
\section{Experimental Results}
\label{sec:results}

\subsection{Experimental Hardware and Software Environment}
\label{sec:platform}

{The physical architecture operates via an ESP32 microcontroller acquiring ultrasonic data and issuing 5~kHz PWM signals (8-bit resolution). The computational host executing the observers and IPOPT optimization is an AMD Ryzen~7 7735HS processor (8 cores/16 threads, 31.2~GB RAM, Windows~11) running Python~3.10.11 and CasADi~3.7.2.} The layout of the hydraulic system and the supporting MOSFET circuitry are depicted in Figure~\ref{fig:plant} and Figure~\ref{fig:circuit}, respectively. The main hardware components of the experimental platform, including the microcontroller, transducers, actuators, and drivers, are summarized in Table~\ref{tab:hardware}. 

\begin{figure}[H]
  %\centering
  \includegraphics[width=0.75\linewidth]{img/imagen_actualizada.jpeg}
  \caption{Coupled three-tank experimental system. Tanks~1 and~3 are actively actuated and measured. Tank~2 is unmeasured within the control loop, utilizing a dedicated, independent ultrasonic sensor strictly for post-hoc validation.}
  \label{fig:plant}
\end{figure}

%\supplementarymaterials{%
  Video S1 shows the real-time closed-loop operation of the physical plant: \url{https://www.youtube.com/watch?v=U-SSxNt9jpA} (accessed on 20 November 2026).

\begin{table}[H]
%  \centering
  \caption{Hardware
 specifications of the experimental platform.}
  \label{tab:hardware}
  \setlength{\tabcolsep}{1.02mm}{
  \begin{tabular}{lll}
    \toprule
    \textbf{Component} & \textbf{Model/Value} & \textbf{Function} \\
    \midrule
    Microcontroller & ESP32 dual core, 240~MHz & Sensing, PWM generation, serial I/O \\
    Transducers     & HC-SR04 (median of 7 pulses) & $h_1,h_3$ (loop control), $h_2$ (validation) \\
    Actuators       & R385 diaphragm pumps, 12~V & Fluid recirculation \\
    Drivers         & IRLZ44N + 1N4007, PWM 5~kHz & Low-side switching \\
    Tanks           & Cylindrical PVC, $\emptyset\,6.2$~cm & Process vessels \\
    \bottomrule
  \end{tabular}}
\end{table}

\begin{figure}[H]
%  \centering
  \includegraphics[width=0.85\linewidth]{img/circuit_diagram.png}
  \caption{MOSFET drive circuit per pump channel ($R_g=100~\Omega$, $R_{pd}=10~\mathrm{k}\Omega$, D1~=~1N4007, $C=100~\mu\mathrm{F}$).}
  \label{fig:circuit}
\end{figure}



\subsection{Global Closed-Loop Performance}

{Table~\ref{tab:summary} highlights the aggregated performance across the ten physical runs based on the gains implemented during the live experiments. The Koopman pipeline achieves a $33.7\%$ reduction in controller tracking error (estimate vs.\ reference) and a $19.3\%$ reduction in observation error against ground truth ($1.139 \pm 0.057$ vs.\ $1.412 \pm 0.064$~cm). However, the absolute physical RMSE of the plant against the setpoint ($\mathrm{RMSE}^\mathrm{real}$) shows no statistically detectable difference between the two methods ($0.804\pm0.031$ vs.\ $0.799\pm0.011$~cm, Welch test $p=0.74$). Because the MPC drives the estimated state toward the reference, any observer bias is transferred, with opposite sign, to the physical plant (Section~\ref{sec:per_segment}). A lower estimate-to-reference error therefore does not imply better physical setpoint tracking. The observation RMSE ranges are disjoint across runs ($[1.093,1.231]$ vs.\ $[1.339,1.507]$~cm), as are the estimate-to-reference ranges ($[0.567,0.633]$ vs.\ $[0.812,0.993]$~cm), so the exact Mann--Whitney $p$-value is $0.0079$ in both cases, the smallest attainable for five versus five runs.}

{Run-level uncertainty. Two-sided $95\%$ Student-$t$ confidence intervals (five runs per method, $t_{0.975,4}=2.776$), computed from the per-run values of Table~\ref{tab:comparison_a}, are: observation RMSE $1.139$ [$1.068$, $1.210$]~cm (Koopman) and $1.412$ [$1.333$, $1.491$]~cm (EKF); estimate-to-reference RMSE $0.595$ [$0.560$, $0.630$]~cm and $0.897$ [$0.815$, $0.979$]~cm; physical RMSE $0.804$ [$0.766$, $0.842$]~cm and $0.799$ [$0.785$, $0.813$]~cm. The as-run observation advantage of the Koopman observer is $0.27$~cm ($95\%$ CI [$0.18$, $0.36$], Welch), of the same order as the $3$~mm resolution of the HC-SR04 and, as shown in Section~\ref{sec:tuning_cv}, not preserved once both observers are tuned. The difference in physical RMSE is $-0.005$~cm ($95\%$ CI [$-0.042$, $0.032$]), which includes zero. The ground-truth sensor was characterized only through its steady-state dispersion ($\sigma\approx 0.45$~cm, $0.35$--$0.58$~cm across runs, Section~\ref{sec:ground_truth}); no independent absolute calibration against a reference scale was performed, which is stated as a limitation in Section~\ref{sec:limitations}. The runs were performed sequentially on a single plant, so these intervals describe repeatability on this platform and not a population-level effect.}

\begin{table}[H]
%  \centering
  \caption{{Performance comparison under as-run experimental gains (mean $\pm$ std over five runs per method, unified steady-state window $N_\mathrm{SS}=310$).}}
  \label{tab:summary}
  \begin{adjustbox}{max width=\linewidth}
  \begin{tabular}{lccc}
    \toprule
    \textbf{Metric} & \textbf{Koopman +
 MPC} & \textbf{EKF + MPC} & \textbf{Comparison (EKF vs.\ Koopman)} \\
    \midrule
    $\mathrm{RMSE}^\mathrm{track}$ (est.\ vs.\ ref.) [cm] & $0.595\pm0.028$ & $0.897\pm0.066$ & Koopman $33.7\%$ lower (MWU $p=0.0079$) \\
    $\mathrm{RMSE}^\mathrm{obs}$ (est.\ vs.\ GT) [cm]    & $1.139\pm0.057$ & $1.412\pm0.064$ & Koopman $19.3\%$ lower (MWU $p=0.0079$) \\
    $\mathrm{RMSE}^\mathrm{real}$ (GT vs.\ ref.) [cm]    & $0.804\pm0.031$ & $0.799\pm0.011$ & No detectable difference (Welch $p=0.74$) \\
    $\mathrm{MAE}^\mathrm{track}$ [cm]                   & $0.508\pm0.028$ & $0.770\pm0.072$ & Koopman $34.0\%$ lower \\
    $\max|e^\mathrm{track}|$ [cm]                        & $1.584\pm0.106$ & $2.045\pm0.076$ & Koopman $22.5\%$ lower \\
    {Observer step time, isolated [ms]}  & {$0.025$} & {$0.087$ (AD) / $0.525$ (FD)} & {Both $\ll T_s$ (Section~\ref{sec:timing})} \\
    \bottomrule
  \end{tabular}
  \end{adjustbox}
\end{table}

Tables~\ref{tab:comparison_a} and~\ref{tab:comparison_b} break down the metrics for individual runs. The distributions of the tracking RMSE are disjoint. Moreover, the total control variation (TV) for actuation $u_1$ remains almost identical ($34.27$ vs.\ $34.93$, $p=0.39$), {indicating} that neither method achieves accuracy by over-actuating the physical pumps.

\begin{table}[H]
  %\centering
  \caption{Steady-state error metrics per run (cm, unified window).}
  \label{tab:comparison_a}
  \footnotesize
  \setlength{\tabcolsep}{4.7mm}{
  \begin{tabular}{llcccc}
    \toprule
    \textbf{Method} & \textbf{Run} & \textbf{RMSE\boldmath{$^\mathrm{track}$}} & \textbf{RMSE\boldmath{$^\mathrm{obs}$}} & \textbf{RMSE\boldmath{$^\mathrm{real}$}} & \textbf{MAE\boldmath{$^\mathrm{track}$}} \\
    \midrule
    Koopman & 1 & 0.568 & 1.093 & 0.799 & 0.484 \\
    Koopman & 2 & 0.598 & 1.106 & 0.769 & 0.517 \\
    Koopman & 3 & 0.610 & 1.231 & 0.853 & 0.515 \\
    Koopman & 4 & 0.567 & 1.109 & 0.797 & 0.477 \\
    Koopman & 5 & 0.633 & 1.159 & 0.803 & 0.546 \\
    \midrule
    EKF & 1 & 0.812 & 1.339 & 0.797 & 0.676 \\
    EKF & 2 & 0.920 & 1.433 & 0.786 & 0.806 \\
    EKF & 3 & 0.993 & 1.507 & 0.798 & 0.868 \\
    EKF & 4 & 0.889 & 1.370 & 0.797 & 0.752 \\
    EKF & 5 & 0.873 & 1.410 & 0.817 & 0.747 \\
    \midrule
    \multicolumn{2}{l}{Koopman
 mean $\pm\sigma$} & $0.595\pm0.028$ & $1.139\pm0.057$ & $0.804\pm0.031$ & $0.508\pm0.028$ \\
    \multicolumn{2}{l}{EKF mean $\pm\sigma$} & $0.897\pm0.066$ & $1.412\pm0.064$ & $0.799\pm0.011$ & $0.770\pm0.072$ \\
    \bottomrule
  \end{tabular}}
\end{table}

\begin{table}[H]
%  \centering
  \caption{Tracking-error bias, dispersion, and total actuation variation (TV) per run.}
  \label{tab:comparison_b}
  \footnotesize
\setlength{\tabcolsep}{2.9mm}{
  \begin{tabular}{llccccc}
    \toprule
    \textbf{Method} & \textbf{Run} & \textbf{Bias\boldmath{$^\mathrm{track}$}} & \boldmath{$\sigma_\mathrm{err}$} & \textbf{Max err.} & \textbf{TV} \boldmath{$u_1$} & \textbf{TV} \boldmath{$u_3$} \\
    \midrule
    Koopman & 1 & $-0.453$ & 0.344 & 1.472 & 35.30 & 4.87 \\
    Koopman & 2 & $-0.483$ & 0.354 & 1.576 & 32.38 & 4.79 \\
    Koopman & 3 & $-0.500$ & 0.349 & 1.760 & 35.16 & 5.09 \\
    Koopman & 4 & $-0.428$ & 0.373 & 1.563 & 35.31 & 4.92 \\
    Koopman & 5 & $-0.508$ & 0.379 & 1.549 & 33.23 & 5.01 \\
    \midrule
    EKF & 1 & $-0.618$ & 0.527 & 2.103 & 35.87 & 5.50 \\
    EKF & 2 & $-0.771$ & 0.502 & 2.055 & 34.57 & 5.53 \\
    EKF & 3 & $-0.846$ & 0.521 & 2.043 & 34.12 & 5.36 \\
    EKF & 4 & $-0.664$ & 0.591 & 2.106 & 35.60 & 5.31 \\
    EKF & 5 & $-0.701$ & 0.521 & 1.920 & 34.48 & 5.54 \\
    \midrule
    \multicolumn{2}{l}{Koopman mean 
 $\pm\sigma$} & $-0.474\pm0.033$ & $0.360\pm0.015$ & $1.584\pm0.106$ & $34.27\pm1.38$ & $4.94\pm0.12$ \\
    \multicolumn{2}{l}{EKF mean $\pm\sigma$} & $-0.720\pm0.090$ & $0.533\pm0.034$ & $2.045\pm0.076$ & $34.93\pm0.76$ & $5.45\pm0.11$ \\
    \bottomrule
  \end{tabular}}
\end{table}

\subsection{Individual Closed-Loop Runs and Trajectory Overlay}

Figures~\ref{fig:koopman_result}--\ref{fig:overlay} graphically illustrate the dynamic performance of the run presenting an observation RMSE closest to the method median (Koopman run~4 and EKF run~5).

\begin{figure}[H]
  %\centering
  \includegraphics[width=11.2 cm]{img/koopman_best_run.png}
  \caption{Dynamic physical response of Koopman + MPC (run~4, median observation RMSE). The panels display measured levels $h_1$ and $h_3$, unmeasured estimate $\hat h_2$ against ground truth, strict tracking errors (with shaded transient exclusion windows), and requisite pump control signals. {The RMSE annotations in the title bar were computed with the unified steady-state window of Section~\ref{sec:protocol} and match the run~4 row of Table~\ref{tab:comparison_a}. The bottom-right panel of this figure reports the total step interval ($\approx 2.1$~s, dominated by the ESP32 serial transmission latency of Section~\ref{sec:ground_truth}), not the algorithmic cost, which is reported separately in Section~\ref{sec:timing}.}}
  \label{fig:koopman_result}
\end{figure}
\unskip
\begin{figure}[H]
 % \centering
  \includegraphics[width=11.2 cm]{img/ekf_best_run.png}
  \caption{Dynamic physical response of EKF + MPC (run~5, median observation RMSE). Layout matches Figure~\ref{fig:koopman_result}. {The RMSE annotations in the title bar were computed with the unified steady-state window of Section~\ref{sec:protocol} and match the run~5 row of Table~\ref{tab:comparison_a}. The bottom-right panel of this figure reports the total step interval ($\approx 2.1$~s, dominated by the ESP32 serial transmission latency of Section~\ref{sec:ground_truth}), not the algorithmic cost, which is reported separately in Section~\ref{sec:timing}.}}
  \label{fig:ekf_result}
\end{figure}
\unskip
\begin{figure}[H]
%  \centering
  \includegraphics[width=1\linewidth]{img/overlay_best_runs.png}
  \caption{Direct overlay of the median runs: Koopman (run~4, blue dashed) and EKF (run~5, red dotted), each compared against its own ground-truth sensor trace. {Each observation RMSE in the legend is computed against the ground truth of that run's own experiment (Koopman run~4: $\approx 1.109$~cm; EKF run~5: $\approx 1.410$~cm), matching the corresponding rows of Table~\ref{tab:comparison_a}. Because the two experiments were performed sequentially, the two ground-truth traces differ slightly and are both shown for transparency.}}
  \label{fig:overlay}
\end{figure}

\subsection{Per-Segment Decomposition and Physical Tracking Consistency}
\label{sec:per_segment}
The aggregate metrics reported above pool the three operating setpoints of the reference profile. Because the nonlinear Torricelli dynamics and the coupling flows change with the operating level, the accuracy of both observers may also vary across setpoints. To verify that the observed trends are not driven by a single operating condition, the steady-state RMSE metrics were recomputed separately for each segment ($0.10$, $0.20$, and $0.15$~m), using the same unified steady-state window.

{Table~\ref{tab:seg_summary} breaks down the RMSE metrics segment-by-segment. The observation error favors the Koopman observer across all regimes ($1.412$ vs.\ $1.860$~cm in Seg.~1; $1.274$ vs.\ $1.394$~cm in Seg.~2; $0.556$ vs.\ $0.793$~cm in Seg.~3). However, evaluating the physical RMSE against ground truth shows no statistically detectable difference at any operating setpoint (Welch $p=0.48$, $0.37$, $0.18$; Table~\ref{tab:seg_summary}).}

\begin{table}[H]
  %\centering
  \caption{Per-segment RMSE (cm, mean $\pm\sigma$ over five runs; $p$-values quoted in the text are Welch tests on the per-run RMSE distributions).}
  \label{tab:seg_summary}
  \resizebox{\linewidth}{!}{%
  \begin{tabular}{lcccc}
    \toprule
    \textbf{Segment} & \textbf{Physical (Koopman)} & \textbf{Physical (EKF)} & \textbf{Obs.\ (Koopman)} & \textbf{Obs.\ (EKF)} \\
    \midrule
    Seg.~1 ($0.10$~m) & $0.942\pm0.077$ & $0.973\pm0.053$ & $1.412\pm0.087$ & $1.860\pm0.082$ \\
    Seg.~2 ($0.20$~m) & $0.927\pm0.062$ & $0.895\pm0.044$ & $1.274\pm0.090$ & $1.394\pm0.056$ \\
    Seg.~3 ($0.15$~m) & $0.452\pm0.019$ & $0.425\pm0.035$ & $0.556\pm0.047$ & $0.793\pm0.066$ \\
    \bottomrule
  \end{tabular}}
\end{table}

{Table~\ref{tab:real_bias} leverages the algebraic identity $e^\mathrm{real}_k = e^\mathrm{track}_k - e^\mathrm{obs}_k$ to map structural biases. In Segments~1 and~2, both estimators underestimate the actual fluid level by $1.1$ to $1.8$~cm. Attempting to close the control error, the MPC overdrives the physical plant $+0.7$ to $+0.9$~cm above the commanded $0.10$ and $0.20$~m setpoints. In magnitude, the EKF has a smaller $|$physical bias$|$ than the Koopman observer by $0.58$~mm in Segment~2 ($+0.708$ vs.\ $+0.766$~cm) and by $0.63$~mm in Segment~3 ($-0.121$ vs.\ $-0.184$~cm), while the Koopman observer is smaller by $0.42$~mm in Segment~1 ($+0.830$ vs.\ $+0.872$~cm). All three differences are below the $3$~mm sensor resolution, so neither observer demonstrates a systematic physical advantage over the other, which confirms that an observer's proximity to the reference state does not nullify static control biases.}

\begin{table}[H]
 % \centering
  \caption{{Per-segment steady-state bias decomposition (cm, mean over five runs).}}
  \label{tab:real_bias}
  \footnotesize
  \resizebox{\linewidth}{!}{%
  \begin{tabular}{llccc}
    \toprule
    \textbf{Method} & \textbf{Error Metric} & \textbf{Seg.~1 (}\boldmath{$0.10$}~\textbf{m)} & \textbf{Seg.~2 (}\boldmath{$0.20$}~\textbf{m)} & \textbf{Seg.~3 (}\boldmath{$0.15$}~\textbf{m)} \\
    \midrule
    Koopman & Tracking: $\hat h_2-h_2^\mathrm{ref}$ & $-0.514$ & $-0.376$ & $-0.531$ \\
            & Observation: $\hat h_2-h_2^\mathrm{real}$ & $-1.344$ & $-1.142$ & $-0.347$ \\
            & Physical: $h_2^\mathrm{real}-h_2^\mathrm{ref}$ & $+0.830$ & $+0.766$ & $-0.184$ \\
    \midrule
    EKF     & Tracking: $\hat h_2-h_2^\mathrm{ref}$ & $-0.918$ & $-0.536$ & $-0.706$ \\
            & Observation: $\hat h_2-h_2^\mathrm{real}$ & $-1.790$ & $-1.244$ & $-0.585$ \\
            & Physical: $h_2^\mathrm{real}-h_2^\mathrm{ref}$ & $+0.872$ & $+0.708$ & $-0.121$ \\
    \bottomrule
  \end{tabular}}
\end{table}

\subsection{Computational Profiling and Hardware Constraints}
\label{sec:timing}

{To quantify the computational load of both observers independently of the telemetry latency, the algorithms were profiled offline with \texttt{perf\_counter\_ns} over 5,000 iterations on the AMD Ryzen~7 host PC (Table~\ref{tab:timing}). The Koopman observer matrix update costs $0.025$~ms. The full EKF recursion, evaluated with finite differences exactly as deployed in the live runs, consumes $0.525$~ms. As an efficient reference, the exact continuous-time Jacobian was also generated with CasADi algorithmic differentiation (AD). The $0.038$~ms figure corresponds exclusively to this Jacobian evaluation (including the Python call overhead), averaged over 5,000 iterations after a 50-iteration warm-up; it does not include the RK4 prediction nor the covariance update. For completeness, the full EKF recursion is $\approx 0.087$~ms with the AD Jacobian and $0.525$~ms with the finite-difference Jacobian used in the live runs.}

{The non-convex MPC optimization requires $14.50$~ms (averaging $6.8$ IPOPT iterations), dominating the algorithmic footprint of both pipelines. The complete Koopman observer-plus-controller pipeline executes in $14.53$~ms on average (P99 $19.96$~ms), corresponding to $0.7\%$ of the $2.1$~s sampling interval; the corresponding EKF pipeline is $15.03$~ms with the finite-difference Jacobian used in the live runs, or $14.59$~ms with the algorithmic-differentiation Jacobian. The $2.1$~s step times recorded during physical experiments were artifacts of the ESP32 transmission protocol blocking the Python serial buffer, not computational bottlenecks. No algorithmic deadlines were breached, so both observers operate well within the real-time budget.}

\begin{table}[H]
%  \centering
  \caption{{Isolated algorithmic execution profiling (AMD Ryzen~7 7735HS, Windows~11, Python~3.10.11, CasADi~3.7.2).}}
  \label{tab:timing}
  \resizebox{\linewidth}{!}{%
  \begin{tabular}{lccc}
    \toprule
    \textbf{Component} & \textbf{Mean (ms)} & \textbf{Std (ms)} & \textbf{P99 (ms)} \\
    \midrule
    Koopman observer update                                  & 0.025 & 0.016 & 0.063 \\
    EKF Jacobian via Algorithmic Differentiation (CasADi)    & 0.038 & 0.027 & 0.123 \\
    EKF full recursion (AD Jacobian)                         & 0.087 & 0.041 & 0.211 \\
    EKF recursion (finite-difference Jacobian, as run)       & 0.525 & 0.173 & 1.116 \\
    MPC Nonlinear Program (IPOPT, $N=20$)                    & 14.50 & 2.043 & 19.90 \\
    \midrule
  Total Algorithmic Pipeline (Koopman + MPC) 
      & 14.53 & 2.04 & 19.96 \\
    \bottomrule
  \end{tabular}}
  
  \noindent{\small{Note: the total P99 of $19.96$~ms is the sum of the per-stage P99 values reported in the table, not the P99 of an end-to-end measurement; the actual end-to-end P99 may be slightly larger due to scheduling jitter, but remains well below the $2.1$~s sampling interval.}}
\end{table}

%\textcolor{blue}{\textit{Note:} the total P99 of $19.96$~ms is the sum of the per-stage P99 values reported in the table, not the P99 of an end-to-end measurement; the actual end-to-end P99 may be slightly larger due to scheduling jitter, but remains well below the $2.1$~s sampling interval.}

\subsection{Transient Response}
\label{sec:transient}

Transient responsiveness evaluated on the ground-truth sensor is quantified by the $90\%$ rise time ($t_{90}$) and \rev{the settling time ($t_s$) into} an absolute $\pm1$~cm band (Table~\ref{tab:transient}). The rise times are comparable: $24.0$ vs.\ $26.4$~s in Segment~1, and $39.6$ vs.\ $43.2$~s in Segment~3. Overshoots remain between $0.42$ and $0.72$~cm ({4.2\% to 10.8\%} 
of the active step magnitude). {One Koopman run in Segment~2 
exhibits a delayed settling time; this run dominates the settling-time standard 
deviation ($65.6 \pm 62.1$~s) and is retained in the statistics for transparency. 
Because part of its transient extends into the retained steady-state window, 
it also inflates the Segment~2 physical and observation RMSE.}

\begin{table}[H]
  %\centering
  \caption{Transient metrics evaluated directly on the ground-truth sensor (mean $\pm\sigma$, five runs).}
  \label{tab:transient}
  \footnotesize
  \begin{adjustbox}{max width=\linewidth}
  \setlength{\tabcolsep}{5.6mm}{
  \begin{tabular}{llccc}
    \toprule
    \textbf{Metric} & \textbf{Method} & \textbf{Seg.~1 (}\boldmath{$0.10$}~\textbf{m)} & \textbf{Seg.~2 (}\boldmath{$0.20$}~\textbf{m)} & \textbf{Seg.~3 (}\boldmath{$0.15$}~\textbf{m)} \\
    \midrule
    $t_{90}$ [s] & Koopman & $24.0\pm1.4$ & $33.6\pm3.3$ & $39.6\pm2.2$ \\
                 & EKF     & $26.4\pm0.9$ & $37.6\pm1.7$ & $43.2\pm5.2$ \\
    \midrule
    \rev{$t_s(\pm1\,\mathrm{cm})$} [s] & Koopman & $24.0\pm1.4$ & $65.6\pm62.1$ & $38.0\pm2.8$ \\
                 & EKF     & $26.0\pm0.0$ & $37.2\pm1.1$ & $41.2\pm4.1$ \\
    \midrule
    Overshoot [cm] & Koopman & $0.57\pm0.23$ & $0.72\pm0.25$ & $0.54\pm0.18$ \\
                   & EKF     & $0.42\pm0.07$ & $0.64\pm0.15$ & $0.45\pm0.11$ \\
    \bottomrule
  \end{tabular}}
  \end{adjustbox}
\end{table}

\subsection{Predictor Baseline and Active Error Rejection}
\label{sec:predictor}

{To demonstrate that output injection actively drives measurement correction, logged telemetry generated across all ten runs was replayed offline through the pure lifted predictor formulation ($L=0$). Because $\rho(\mathbf A)=0.99990<1$, the predictor is mathematically stable, but its time constant $\tau_\mathrm{predictor}=-T_s/\ln\rho(\mathbf A)\approx 2.0\times 10^4$~s is about 25 times larger than the $800$~s horizon. The pure predictor drifts to RMS errors of $5.5$--$6.1$~cm in the unified window (Table~\ref{tab:predictor}); since the predictor is stable, this drift is attributed to model--plant mismatch accumulated over the horizon and not to mathematical instability. In contrast, the full Koopman observer ($L>0$) with the closed-loop matrix $\mathbf A_c=(I-LC)\mathbf A$ ($\rho(\mathbf A_c)=0.93515$, $\tau_e=29.83$~s) keeps the observation RMSE at $1.15$~cm.}

{To stress-test the architecture, an initial estimation error of $+2.0$~cm was injected in $\hat h_2(0)$. The effect of the injection is isolated as the difference $\Delta\hat h_2(t)$ between the replay with injection and the injection-free replay of the same telemetry $(y_k,u_k)$; the offset is considered rejected at the first time after which $|\Delta\hat h_2(t)|\le 0.2$~cm ($10\%$ of the injected value) for all later times. The replay with the logged telemetry gives $7\pm1$~s averaged over the ten runs (range $6$--$8$~s, with a resolution of one sample, $2$~s). As a cross-check, noise-free propagation of the lifted error $e_k=\mathbf A_c^k e_0$ with the identified $\mathbf A_c$ gives $12$~s for the lifted state $|e_{2,k}|$ to fall to $10\%$ of its initial value at the three operating points ($h_2=0.10$, $0.15$, $0.20$~m). The two values differ because the physical injection $|\Delta\hat h_2|=|2\hat z_2\,\Delta\hat z_2+\Delta\hat z_2^2|$ is amplified by $2\hat z_2$, which varies during the transient; the replay criterion on the physical value is met earlier when $\hat z_2$ is small. A modal decomposition of the injected offset onto the eigenmodes of $\mathbf A_c$ explains why the transient is much faster than the slowest mode: the $+2$~cm offset projects only $1.4$--$5.2\%$ onto the slow mode ($|\lambda|=0.9351$, $\tau=29.83$~s) and $52$--$57\%$ onto the second mode ($|\lambda|=0.6565$, $\tau=4.8$~s). The asymptotic constant $\tau_e=29.83$~s therefore characterizes the contraction factor of Proposition~\ref{prop:stability}, not the transient decay time of this particular offset.}

{The unified-window RMSE is practically identical ($1.15$~cm) for the three initializations tested (Figure~\ref{fig:predictor_baseline}). The observer thus settles at a residual floor bounded by $(\|I-LC\|_2\bar r+\|L\|_2\bar v)\sqrt{\kappa(W)}/(1-\gamma)$ rather than converging to the true state: on the real plant it remains at the offset induced by the EDMD model mismatch and the HC-SR04 noise (Table~\ref{tab:real_bias}).}

\begin{table}[H]
 % \centering
  \caption{{Observation RMSE [cm] isolating the active feedback effect (offline telemetry replay).}}
  \label{tab:predictor}
  \footnotesize
  \resizebox{\linewidth}{!}{%
  \begin{tabular}{lcccc}
    \toprule
    & \multicolumn{2}{c}{\textbf{Unified Window}} & \multicolumn{2}{c}{\textbf{Whole Run}} \\
    \cmidrule(lr){2-3}\cmidrule(lr){4-5}
    \textbf{Initial Estimate} \boldmath{$\hat h_2(0)$} & \boldmath{$L>0$} \textbf{(Active Observer)} & \boldmath{$L=0$} \textbf{(Pure Predictor)} & \boldmath{$L>0$} & \boldmath{$L=0$} \\
    \midrule
    $h_2^\mathrm{ref}(0)$ (Experimental config) & $1.151\pm0.048$ & $6.149\pm1.359$ & $1.337\pm0.038$ & $6.035\pm1.333$ \\
    \rev{$h_2^\mathrm{real}(0)$}$+2$~cm          & $1.151\pm0.048$ & $5.708\pm1.303$ & $1.247\pm0.040$ & $5.605\pm1.266$ \\
    \rev{$h_2^\mathrm{real}(0)$}                 & $1.152\pm0.048$ & $5.530\pm1.266$ & $1.244\pm0.041$ & $5.452\pm1.221$ \\
    \bottomrule
  \end{tabular}}
\end{table}

\begin{figure}[H]
  %\centering
  \includegraphics[width=\linewidth]{img/predictor_L0_validation.png}
  \caption{{Replay validation isolating active output injection. The pure predictor ($L=0$, red dotted) drifts due to model--plant mismatch. The active Koopman observer ($L>0$, blue dashed) rejects a $+2.0$~cm initialization offset within a few samples (Section~\ref{sec:predictor}) and tracks the ground truth (green solid) down to a residual floor set by the EDMD model mismatch and the HC-SR04 noise; the steady-state observation bias is $\sim 1.1$--$1.3$~cm \rev{in Segments~1 and~2} (Table~\ref{tab:real_bias}), larger than the $\approx 0.45$~cm sensor noise, indicating that model--plant mismatch dominates the residual.}}
  \label{fig:predictor_baseline}
\end{figure}

%% =====================================================================
\section{Discussion}
\label{sec:discussion}

\subsection{Offline Cross-Validation and Methodological Fairness}
\label{sec:tuning_cv}

{To isolate estimation accuracy from closed-loop interaction, the recorded telemetry $(y_k,u_k)$ of the ten runs was replayed offline through both observers with identical preprocessing and the unified evaluation window. Both observers are deterministic, so each one replayed on its own five logs with the as-run gains reproduces its live trajectories. For the EKF this holds exactly ($1.412\pm0.064$~cm in live and replay, identical run by run: $1.339$, $1.433$, $1.507$, $1.370$, $1.410$~cm). For the Koopman observer, replay on its own five logs gives $1.139\pm0.057$~cm, identical to the live value, as expected for a deterministic linear observer replayed with the same initial condition. Applying each observer to all ten logs (Table~\ref{tab:ablation}) gives the pooled replay values used as the reference for the tuning gains below.}

\begin{table}[H]
 %\centering
  \caption{Replay-based ablation of the observation RMSE against the ground-truth sensor (unified window). ``Own'' denotes the five logs recorded with that observer in closed loop; ``All'' denotes all ten logs. LOO-tuned values are evaluated on ten held-out runs; the percentage improvement is measured against the corresponding live as-run value of each observer ($1.139$~cm for Koopman, $1.412$~cm for EKF).}
  \label{tab:ablation}
  \begin{adjustbox}{max width=\linewidth}
  \setlength{\tabcolsep}{4mm}{
  \begin{tabular}{llccc}
    \toprule
    \textbf{Observer} & \textbf{Configuration} & \textbf{Logs} & \textbf{RMSE\boldmath{$^\mathrm{obs}$} [cm]} & \textbf{vs.\ Live as-Run} \\
    \midrule
    Koopman & as-run, live             & Own (5)  & $1.139\pm0.057$ & --- \\
    Koopman & as-run, replay           & Own (5)  & $1.139\pm0.057$ & $0.0\%$ \\
    Koopman & as-run, replay           & All (10) & $1.151\pm0.048$ & $+1.0\%$ \\
    Koopman & LOO-tuned                & All (10) & $0.655\pm0.023$ & $-43\%$ \\
    \midrule
    EKF     & as-run, live             & Own (5)  & $1.412\pm0.064$ & --- \\
    EKF     & as-run, replay           & Own (5)  & $1.412\pm0.064$ & $0.0\%$ \\
    EKF     & LOO-tuned                & All (10) & $0.600\pm0.020$ & $-58\%$ \\
    \bottomrule
  \end{tabular}}
  \end{adjustbox}
\end{table}

{Under leave-one-run-out cross-validation against the ground-truth sensor (one held-out run at a time, ten folds, all ten logs pooled), both estimators were optimized: $30$ configurations for the Koopman observer and $72$ covariance permutations for the EKF. The EKF tuned to ($\sigma_{h_1h_3}=0.5$, $\sigma_{h_2}=1.0$, $\sigma_R=4.0$~mm) reaches $0.600\pm0.020$~cm; the Koopman observer tuned to ($\alpha_\mathrm{med}=10^{-3}$, $\alpha_\mathrm{un}=0.1$) reaches $0.655\pm0.023$~cm. The comparison that bears on the cause of the as-run gap is the following: the tuned EKF ($0.600$~cm) is more accurate than the as-run Koopman observer ($1.151$~cm) by about $5.5$~mm, while the two tuned observers differ by only $0.55$~mm, with the EKF slightly better (Wilcoxon $p=0.00195$, ten paired folds), less than one fifth of the $3$~mm resolution of the HC-SR04. The as-run observation advantage of the Koopman observer is therefore largely attributable to the EKF covariance configuration deployed in the live runs ($\sigma_{h_2}=2.0$~mm and $\sigma_R=1.5$~mm, versus $1.0$~mm and $4.0$~mm after tuning; the empirical sensor noise floor is $\approx 4.5$~mm). This analysis does not isolate which covariance term is responsible, and it does not by itself exclude a structural difference between the estimators. Since the as-run Koopman gain was not tuned either (it improves by $43\%$), this analysis does not rank the two structures; it shows that the live gap does not persist under tuning. The ground-truth sensor itself has a noise standard deviation of $\approx 0.45$~cm, of the same order as the tuned RMSE, so both observers operate close to the measurement limit of the validation. We emphasize three caveats: (i) the replay does not fully separate model error from tuning, since both models derive from the same calibrated ODE; (ii) the cross-validation budget is asymmetric ($30$ vs.\ $72$ configurations); and (iii) the tuning targets the ground-truth sensor, which is available only offline and is not part of the deployable pipeline.}

\subsection{Algebraic Consistency on the Lifted Space}

{The Koopman observer inherently evolves within unconstrained $\mathbb R^8$ coordinate geometry. This operational freedom permits the lifted state $\hat z_k$ to deviate from the pure algebraic manifold $\{z=\Psi(x)\}$. Specifically, the RMS algebraic residual $z_4-z_1^2z_2^2$ across the experimental telemetry is $3.30\times10^{-2}$~m$^2$, i.e., approximately $125\%$ of the RMS value of $z_4$ itself ($2.65\times10^{-2}$~m$^2$): the lifted estimate does not satisfy the algebraic constraint $z_4=z_1^2z_2^2$ in an RMS sense, and the dependent observables $z_4,z_5,z_6$ evolve under the linear propagation independently of the products they nominally represent. If the physical state is re-lifted at every time step ($\hat z_k\leftarrow\Psi(\hat x_k)$) to force algebraic consistency, the observation RMSE degrades by $+10.3\%$. The observer therefore benefits from the flexibility afforded by unconstrained linear propagation in $\mathbb R^8$ to absorb transient dynamic mismatches.}

\subsection{Study Limitations}
\label{sec:limitations}

{These findings are subject to several limitations. First, the mathematical plant model assumes a pure $T_s=2.0$~s period void of latency, whereas the physical ESP32 acquisition chain operates at roughly $2.10$~s and injects approximately $0.7$~s of actuation delay. This unmodeled delay directly produces a steady-state tracking drift common to both estimation pipelines. Second, the ground-truth HC-SR04 transducer operates with the same noise profile as the control loop sensors ($\approx 0.45$~cm standard deviation from pump jet surface ripples). Third, identifying the EDMD lifted model relied exclusively on simulated data rather than live hardware telemetry, and the ten physical runs were performed on a single plant with a single reference sequence ($0.10\to0.20\to0.15$~m) and no external flow disturbances; testing under broader excitation envelopes remains open. The physical runs are further limited to $800$~s; longer physical runs and additional tests with unseen reference profiles are left for future work. No independent calibration of the validation sensor against an external reference scale was performed; the calibration of the ground-truth sensor rests on its steady-state dispersion and quantization resolution, and any absolute offset remains intertwined with the observer bias. \rev{It is important to delimit how this $\approx 0.45$~cm dispersion affects the reported accuracy. If the sensor deviation is zero-mean and uncorrelated with the estimation error, it adds in quadrature to the measured error, $(\mathrm{RMSE}^\mathrm{obs}_\mathrm{measured})^2 \approx (\mathrm{RMSE}^\mathrm{obs}_\mathrm{true})^2 + \sigma^2$, so the reported $\approx 1.1$--$1.4$~cm as-run values (and the $\approx 0.6$~cm tuned values) should be read as sensor-noise-limited upper bounds rather than exact estimation errors. Because this term is the same for both observers, it does not change their ordering. In the offline replay of Section~\ref{sec:tuning_cv}, both observers are evaluated against identical ground-truth samples, so sample-wise sensor fluctuations are common to both and are accounted for by the paired comparison. In the live comparison of Section~\ref{sec:results}, the runs were performed sequentially, so these fluctuations enter as run-to-run variability and are reflected in the reported confidence intervals. A constant calibration offset, by contrast, does not cancel, as noted above. Once both observers are tuned, their $0.55$~mm gap, although detectable in the paired test, is well below the $3$~mm sensor resolution and is not practically meaningful.} Finally, the ultimate-bound criterion of Proposition~\ref{prop:stability} is conservative, enforcing mathematical boundedness without offering tight thresholds for operational deployment.}

\rev{Because the validation rests on five runs per observer performed sequentially on a single physical setup with a single reference profile, the reported differences should be read as evidence of a repeatable behavior on this specific platform rather than as a population-level or cross-platform claim. Establishing general applicability to other plants and operating envelopes would require multiple setups, more runs per condition, and adequately powered hypothesis testing; the conclusions are scoped accordingly.}

%% =====================================================================
\section{Conclusions}
\label{sec:conclusions}

A Koopman-based observer, designed offline in an eight-observable lifted space, was compared with an EKF on a physical three-tank plant using an independent ground-truth sensor. {With the gains used in the live experiments, the Koopman observer estimated the unmeasured level $19.3\%$ more accurately; however, the physical tracking error of the plant showed no statistically detectable difference between the two pipelines. {A predictor-only baseline confirmed that the output injection actively rejects a deliberately injected $+2.0$~cm offset, decaying it to $10\%$ within approximately $7$~s in the replay ($12$~s under noise-free propagation), while the pure predictor drifts.} Offline leave-one-run-out cross-validation replay subsequently indicated that tuning the EKF closes and slightly inverts this gap: the EKF reaches $0.600$~cm and the Koopman observer $0.655$~cm RMSE, with the EKF marginally better (Wilcoxon $p=0.00195$ over ten paired folds), and both algorithms achieve comparable estimation precision, limited by the $\approx 0.45$~cm noise of the validation sensor; the $0.55$~mm residual difference is far below the sensor resolution. Isolated algorithmic profiling confirmed that the full recursion of an EKF with an algorithmic-differentiation Jacobian ($0.087$~ms) is comparable to the linear Koopman update ($0.025$~ms); the EKF Jacobian evaluation alone requires $0.038$~ms. Both are negligible against the non-convex model predictive optimization, which alone requires $14.50$~ms on average (P99 $19.90$~ms), with a complete Koopman observer-plus-controller pipeline of $14.53$~ms on average (P99 $19.96$~ms).}

\rev{We therefore refrain from claiming uniform superiority of the proposed observer. The as-run observation-accuracy advantage reflects the specific fixed gains used in the live experiments, and the cross-validated replay shows that it largely vanishes once both observers are tuned (the EKF is in fact marginally better, by $0.55$~mm, below the $3$~mm sensor resolution). The accurate statement is that the Koopman observer is a stable, computationally efficient, and competitive alternative to the EKF, matching its estimation accuracy after tuning while providing a linear time-invariant structure and an a~priori boundedness certificate, rather than an estimator that outperforms the EKF.}

{The contribution of the Koopman formulation resides in its mathematical structure rather than in computational or operational superiority. It delivers EKF-grade estimation accuracy while relying on a fully linear, time-invariant matrix synthesized offline. It requires no real-time Jacobian linearizations, avoids local Taylor series expansion, and provides an ultimate-boundedness framework for the lifted estimation error of the identified model. Future work will embed the full observer--MPC algorithm on the microcontroller, using an embedded NLP/SQP solver (or a convexified reformulation of the cost) to eliminate the telemetry latency.}

\rev{Regarding extensibility, the methodology carries over to larger cascaded hydraulic networks ($n$-tank systems, water-distribution branches) and to other nonlinear industrial processes with a dominant, physically interpretable nonlinearity---for instance continuous stirred-tank reactors (Arrhenius kinetics), heat exchangers, or pneumatic actuators---provided the dictionary is enlarged with the corresponding physics-informed observables and enough excitation data are available to identify an accurate EDMD model within the operating envelope. The practical limit is the growth of the lifted dimension $n_z$ and of the data requirement with system size; since the observer update scales as $O(n_z^2)$ while the EKF Jacobian scales with the physical dimension, the structural advantages (no online linearization, a~priori boundedness) are expected to be most valuable at moderate $n_z$ and to erode if an accurate predictor demands a very large dictionary or if strong unmodeled regime changes fall outside the training envelope, where online EDMD adaptation would be required.}
\vspace{6pt}

%% =====================================================================
\supplementary{The following supporting information can be downloaded at \linksupplementary{s1}: Video S1: Real-time closed-loop operation of the physical plant.}

\authorcontributions{%
  W.S.H.: Conceptualization, Software, Experimental design and execution,
  Formal analysis, Writing---original draft.
  C.A.T.-P.: Methodology, Software, Data curation, \mbox{Writing---review and editing}.
  D.T.-C.: Supervision, Methodology review, Writing---review and editing.
  All authors have read and agreed to the published version of the manuscript.
}
\funding{This research received no external funding.}
\institutionalreview{Not applicable.}
\informedconsent{Not applicable.}
\dataavailability{The experimental logs (CSV), the ESP32 firmware, the acquisition and control scripts as used in the experiments, the offline analysis scripts and the software environment (package versions) are available at
 \url{https://github.com/wshernandezm/koopman-three-tank-benchmark} (accessed on 20 November 2026).}
\conflictsofinterest{The authors declare no conflicts of interest.}

%% =====================================================================
\appendixtitles{yes}
\appendixstart
\appendix

\section{Identified Model and Observer System Matrices}
\label{app:matrices}

{%
The discrete-time state matrix $\mathbf{A} \in \mathbb{R}^{8 \times 8}$ identified via regularized EDMD and spectrally stabilized ($\delta = 10^{-4}$), alongside the input matrix $\mathbf{B} \in \mathbb{R}^{8 \times 2}$ and the time-invariant observer gain $L \in \mathbb{R}^{8 \times 4}$ synthesized via DARE in predictor form $L=\mathbf AXC^\top(CXC^\top+R)^{-1}$ ($R=I_4$, $Q=I_8$) for the physical experiments. The closed-loop error matrix is $\mathbf A_c=(I-LC)\mathbf A$. The identified matrices satisfy $\rho(\mathbf{A})=0.99990$, $\rho(\mathbf A_c)=0.93515$, and maximum gain magnitude $|L|_{\max} = 0.5867$.

\vspace{0.4cm}
\begin{equation}
\resizebox{0.85\linewidth}{!}{$
\mathbf{A} = \begin{bmatrix}
  0.7941 & 0.1507 & 0.0724 & 0.0216 & 0.0302 & 0.0344 & 0.0143 & -0.0942 \\
  0.1946 & 0.7226 & 0.0680 & 0.0548 & 0.0739 & -0.1497 & 0.0084 & 0.0127 \\
  0.0365 & 0.0833 & 0.8535 & 0.0173 & -0.0707 & -0.0142 & -0.0340 & 0.0932 \\
  -0.0516 & 0.0113 & 0.0399 & 0.5721 & 0.3169 & 0.1792 & 0.1078 & -0.1236 \\
  0.0516 & 0.0042 & -0.0619 & 0.3049 & 0.3937 & 0.2134 & -0.1147 & 0.1401 \\
  0.0352 & -0.0175 & -0.0259 & 0.1833 & 0.1875 & 0.5980 & -0.0418 & 0.0639 \\
  -0.0316 & 0.0987 & -0.0783 & 0.2436 & -0.0136 & -0.2643 & 0.8455 & 0.1724 \\
  -0.0984 & 0.0471 & 0.0607 & -0.0872 & 0.1991 & -0.0421 & 0.1493 & 0.8004
\end{bmatrix}
$}
\end{equation}

\vspace{0.2cm}
\begin{equation}
\resizebox{0.35\linewidth}{!}{$
\mathbf{B} = \begin{bmatrix}
  0.0194 & -0.0005 \\
  0.0023 & 0.0007 \\
  0.0001 & 0.0211 \\
  0.0029 & 0.0002 \\
  0.0003 & 0.0026 \\
  0.0022 & 0.0028 \\
  0.0159 & -0.0001 \\
  0.0001 & 0.0152
\end{bmatrix}
$}
\end{equation}

The output selector matrix $C \in \mathbb{R}^{4 \times 8}$ extracts the measurable combinations associated with $h_1$ and $h_3$:
\begin{equation}
C = \begin{bmatrix}
1 & 0 & 0 & 0 & 0 & 0 & 0 & 0 \\
0 & 0 & 1 & 0 & 0 & 0 & 0 & 0 \\
0 & 0 & 0 & 0 & 0 & 0 & 1 & 0 \\
0 & 0 & 0 & 0 & 0 & 0 & 0 & 1
\end{bmatrix}.
\end{equation}

The experimental observer gain matrix $L \in \mathbb{R}^{8 \times 4}$ synthesized via DARE is:
\begin{equation}
\resizebox{0.6\linewidth}{!}{$
L = \begin{bmatrix}
  0.5225 & 0.0485 & 0.0257 & -0.0442 \\
  0.2277 & 0.1175 & 0.1043 & 0.0318 \\
  0.0295 & 0.5365 & -0.0048 & 0.0543 \\
  0.1977 & -0.1857 & 0.0683 & 0.0538 \\
  0.2178 & -0.2120 & -0.0837 & 0.1917 \\
  0.2155 & -0.2040 & -0.0759 & 0.1508 \\
  -0.0023 & -0.0319 & 0.5867 & 0.1208 \\
  -0.0441 & 0.0348 & 0.1149 & 0.4999
\end{bmatrix}
$}
\end{equation}
}

\begin{adjustwidth}{-\extralength}{0cm}
\reftitle{References}
\begin{thebibliography}{999}

\bibitem[Johansson(2000)]{johansson2000}
Johansson, K.H.
\newblock The quadruple-tank process: A multivariable laboratory process with an adjustable zero.
\newblock {\em IEEE Trans. Control Syst. Technol.} {\bf 2000}, {\em 8},~456--465.
\newblock {\url{https://doi.org/10.1109/87.845876}}.

\bibitem[Simon(2006)]{simon2006}
Simon, D.
\newblock {\em Optimal State Estimation: Kalman, {H}$_\infty$, and Nonlinear
  Approaches}; Wiley: Hoboken, NJ, USA,  2006.


\bibitem[Julier and Uhlmann(2004)]{julier2004}
Julier, S.; Uhlmann, J.
\newblock Unscented filtering and nonlinear estimation.
\newblock {\em Proc. IEEE} {\bf 2004}, {\em 92},~401--422.
\newblock {\url{https://doi.org/10.1109/JPROC.2003.823141}}.

\bibitem[Rao et~al.(2003)Rao, Rawlings, and Mayne]{rao2003}
Rao, C.; Rawlings, J.; Mayne, D.
\newblock Constrained state estimation for nonlinear discrete-time systems:
  Stability and moving horizon approximations.
\newblock {\em IEEE Trans. Autom. Control} {\bf 2003}, {\em 48},~246--258.
\newblock {\url{https://doi.org/10.1109/TAC.2002.808470}}.

\bibitem[Williams et~al.(2015)Williams, Kevrekidis, and Rowley]{williams2015}
Williams, M.; Kevrekidis, I.; Rowley, C.
\newblock A data-driven approximation of the {K}oopman operator: {E}xtending
  dynamic mode decomposition.
\newblock {\em J. Nonlinear Sci.} {\bf 2015}, {\em 25},~1307--1346.
\newblock {\url{https://doi.org/10.1007/s00332-015-9258-5}}.

\bibitem[Korda and Mezi\'{c}(2018)]{korda2018}
Korda, M.; Mezi\'{c}, I.
\newblock Linear predictors for nonlinear dynamical systems: {K}oopman operator meets model predictive control.
\newblock {\em Automatica} {\bf 2018}, {\em 93},~149--160.
\newblock {\url{https://doi.org/10.1016/j.automatica.2018.03.046}}.

\bibitem[Koopman(1931)]{koopman1931}
Koopman, B.O.
\newblock Hamiltonian Systems and Transformation in Hilbert Space.
\newblock {\em Proc. Natl. Acad. Sci. USA} {\bf 1931},
  {\em 17},~315--318.
\newblock {\url{https://doi.org/10.1073/pnas.17.5.315}}.

\bibitem[Mezi\'{c}(2005)]{mezic2005}
Mezi\'{c}, I.
\newblock Spectral properties of dynamical systems, model reduction and
  decompositions.
\newblock {\em Nonlinear Dyn.} {\bf 2005}, {\em 41},~309--325.
\newblock {\url{https://doi.org/10.1007/s11071-005-2824-x}}.

\bibitem[Brunton et~al.(2016)Brunton, Brunton, Proctor, and Kutz]{brunton2016}
Brunton, S.; Brunton, B.; Proctor, J.; Kutz, J.
\newblock {K}oopman invariant subspaces and finite linear representations of
  nonlinear dynamical systems for control.
\newblock {\em PLoS ONE} {\bf 2016}, {\em 11},~e0150171.
\newblock {\url{https://doi.org/10.1371/journal.pone.0150171}}.

\bibitem[Tellez-Castro et~al.(2022)Tellez-Castro, Garcia-Tenorio, Mojica-Nava,
  Sofrony, and Vande~Wouwer]{tellez2022}
Tellez-Castro, D.; Garcia-Tenorio, C.; Mojica-Nava, E.; Sofrony, J.;
  Vande~Wouwer, A.
\newblock Data-driven predictive control of interconnected systems using the
  {K}oopman operator.
\newblock {\em Actuators} {\bf 2022}, {\em 11},~151.
\newblock {\url{https://doi.org/10.3390/act11060151}}.

\bibitem[Hernandez et~al.(2026)Hernandez, Ramirez-Murillo, and
  Tellez-Castro]{hernandez2026}
Hernandez, W.; Ramirez-Murillo, H.; Tellez-Castro, D.
\newblock Data-Driven Koopman-Based Model Predictive Control for a Three-Tank
  Hydraulic System.
\newblock In Proceedings of the 2026 12th International Conference on Control,
  Decision and Information Technologies (CoDIT), Bari, Italy, 13--16 July 2026; pp. 2060--2065.
\newblock {\url{https://doi.org/10.1109/CoDIT70676.2026.11631175}}.

\bibitem[Aponte and Tellez-Castro(2025)]{aponte2025}
Aponte, D.S.; Tellez-Castro, D.
\newblock Koopman-Based MPC for the Quadruple Tank System: An Experimental
  Comparison.
\newblock In Proceedings of the 2025 IEEE 7th Colombian Conference on Automatic
  Control (CCAC), Pereira, Colombia, 14--17 October 2025; pp. 1--6.
\newblock {\url{https://doi.org/10.1109/CCAC64704.2025.11259347}}.

\bibitem[Durán-Siguenza et~al.(2025)Durán-Siguenza, Minchala,
  Garza-Castañón, and Zhang]{DURANSIGUENZA2025}
Durán-Siguenza, J.F.; Minchala, L.I.; Garza-Castañón, L.E.; Zhang, H.
\newblock Control based on the Koopman operator: A comprehensive review.
\newblock {\em J. Frankl. Inst.} {\bf 2025}, {\em
  362},~108256.
\newblock
  {\url{https://doi.org/10.1016/j.jfranklin.2025.108256}}.

\bibitem[Netto and Mili(2018)]{netto2018}
Netto, M.; Mili, L.
\newblock A robust data-driven {K}oopman {K}alman filter for power systems dynamic state estimation.
\newblock {\em IEEE Trans. Power Syst.} {\bf 2018}, {\em 33},~7228--7237.
\newblock {\url{https://doi.org/10.1109/TPWRS.2018.2846744}}.

\bibitem[Surana and Banaszuk(2016)]{surana2016}
Surana, A.; Banaszuk, A.
\newblock Linear observer synthesis for nonlinear systems using {K}oopman
  operator framework.
\newblock {\em IFAC-PapersOnLine} {\bf 2016}, {\em 49},~716--723.
\newblock {\url{https://doi.org/10.1016/j.ifacol.2016.10.250}}.

\bibitem[Surana et~al.(2017)Surana, Williams, Morari, and Banaszuk]{Surana2017}
Surana, A.; Williams, M.O.; Morari, M.; Banaszuk, A.
\newblock Koopman operator framework for constrained state estimation.
\newblock In Proceedings of the 2017 IEEE 56th Annual Conference on Decision
  and Control (CDC),  Melbourne, VIC, Australia, 12--15 December 2017; pp. 94--101.
\newblock {\url{https://doi.org/10.1109/CDC.2017.8263649}}.

\bibitem[Guo et~al.(2026)Guo, Forbes, and Barfoot]{kiloekf2026}
Guo, Z.C.; Forbes, J.R.; Barfoot, T.D.
\newblock KILO-EKF: Koopman-Inspired Learned Observations Extended Kalman
  Filter.  \emph{arXiv} \textbf{2026}, arXiv:2601.12463.

\bibitem[W\"{a}chter and Biegler(2006)]{wachter2006}
W\"{a}chter, A.; Biegler, L.T.
\newblock On the implementation of an interior-point filter line-search algorithm for large-scale nonlinear programming.
\newblock {\em Math. Program.} {\bf 2006}, {\em 106},~25--57.
\newblock {\url{https://doi.org/10.1007/s10107-004-0559-y}}.

\bibitem[Andersson et~al.(2019)Andersson, Gillis, Horn, Rawlings, and
  Diehl]{andersson2019}
Andersson, J.; Gillis, J.; Horn, G.; Rawlings, J.; Diehl, M.
\newblock {CasADi}: A software framework for nonlinear optimization and optimal
  control.
\newblock {\em Math. Program. Comput.} {\bf 2019}, {\em 11},~1--36.
\newblock {\url{https://doi.org/10.1007/s12532-018-0139-4}}.

\end{thebibliography}
\PublishersNote{}
%\bibliography{references}
\end{adjustwidth}
\end{document}
